\documentclass[11pt]{article}
\usepackage[margin=1in]{geometry}
\usepackage{amsmath,amssymb,bm,booktabs,array,longtable}
\usepackage{enumitem}
\usepackage[hidelinks]{hyperref}
\usepackage{natbib}
\newcommand{\R}{\mathbb{R}}
\newcommand{\ones}{\mathbf{1}}
\newcommand{\diag}{\operatorname{diag}}
\newcommand{\col}{\operatorname{col}}
\newcommand{\ScaleBridge}{\textsc{ScaleBridge}}

\title{ScaleBridge Phase E0 E0-3: Current RC Flavours v2}
\author{ScaleBridge / PhD\_Code\_Framework}
\date{Math freeze: 26 August 2026}
\begin{document}
\maketitle

\section{Status and literature relationship}
Low-order grey-box RC families are widely used for building thermal modeling \citep{BacherMadsen2011HeatDynamics,WangChenLi2019RCSimplification}.  Recent comparative work explicitly evaluates the labels 1R1C, 2R2C, 3R2C, and 4R3C \citep{PanagiEtAl2026GreyBoxTuning}; broader recent studies emphasize that topology choice materially affects accuracy, computation, and identifiability \citep{ChenKorolijaRovas2026SingleZoneRC}.  However, RC labels are not universally topology-unique across the literature.  Therefore the exact graphs below are the \ScaleBridge{} ratified built-in topologies; literature citations support the RC-family concept, not an assertion that every cited paper uses these exact graphs.

\section{1R1C canonical}
State: observed zone air $T_a$.  Parameters: $C_a>0$, $R_{ao}>0$.  All applicable atomic thermal-power ports route fully to air:
\[
C_a\dot T_a=\frac{T_o-T_a}{R_{ao}}+Q_{AC}+Q_{ZIC}+Q_{ZIR}+Q_{Sol1}+Q_{Sol2}.
\]
No canonical $\eta$ is introduced because all declared thermal power must remain represented.  Any paper-specific reduced-radiative-participation 1R1C is a separate model variant.

\section{2R2C canonical}
States: observed air $T_a$ and latent internal mass $T_m$. Parameters: $C_a,C_m,R_{ao},R_{am}>0$. Define
\[
Q_c=Q_{AC}+Q_{ZIC}+Q_{Sol1},\qquad Q_r=Q_{ZIR}+Q_{Sol2}.
\]
With $0\le\eta_r\le1$,
\begin{align}
C_a\dot T_a &= \frac{T_o-T_a}{R_{ao}}+\frac{T_m-T_a}{R_{am}}+Q_c+(1-\eta_r)Q_r,\\
C_m\dot T_m &= \frac{T_a-T_m}{R_{am}}+\eta_rQ_r.
\end{align}
Summing gives
\[
C_a\dot T_a+C_m\dot T_m=\frac{T_o-T_a}{R_{ao}}+Q_c+Q_r.
\]

\section{3R2C canonical}
States remain observed air $T_a$ and latent \emph{internal mass} $T_m$.  Add an outdoor--mass resistance $R_{om}>0$:
\begin{align}
C_a\dot T_a &= \frac{T_o-T_a}{R_{ao}}+\frac{T_m-T_a}{R_{am}}+Q_c+(1-\eta_r)Q_r,\\
C_m\dot T_m &= \frac{T_o-T_m}{R_{om}}+\frac{T_a-T_m}{R_{am}}+\eta_rQ_r.
\end{align}
Hence
\[
C_a\dot T_a+C_m\dot T_m=\frac{T_o-T_a}{R_{ao}}+\frac{T_o-T_m}{R_{om}}+Q_c+Q_r.
\]
The exact topology above is a \ScaleBridge{} design choice within the broader low-order RC modeling family.

\section{4R3C canonical}
States: observed air $T_a$, latent envelope $T_e$, latent internal mass $T_m$.  Parameters $C_a,C_e,C_m,R_{ao},R_{ae},R_{eo},R_{am}>0$.  The four resistance edges are outdoor--air, air--envelope, envelope--outdoor, and air--internal-mass.  Convective/direct heat routes to air.  Radiative heat uses
\[
\gamma_r=\begin{bmatrix}\gamma_{a,r}\\\gamma_{e,r}\\\gamma_{m,r}\end{bmatrix},\qquad
\gamma_{a,r},\gamma_{e,r},\gamma_{m,r}\ge0,\qquad
\gamma_{a,r}+\gamma_{e,r}+\gamma_{m,r}=1.
\]
Let $Q_a=Q_c+\gamma_{a,r}Q_r$, $Q_e=\gamma_{e,r}Q_r$, $Q_m=\gamma_{m,r}Q_r$. Then
\begin{align}
C_a\dot T_a &= \frac{T_o-T_a}{R_{ao}}+\frac{T_e-T_a}{R_{ae}}+\frac{T_m-T_a}{R_{am}}+Q_a,\\
C_e\dot T_e &= \frac{T_o-T_e}{R_{eo}}+\frac{T_a-T_e}{R_{ae}}+Q_e,\\
C_m\dot T_m &= \frac{T_a-T_m}{R_{am}}+Q_m.
\end{align}
Summing gives
\[
C_a\dot T_a+C_e\dot T_e+C_m\dot T_m
=\frac{T_o-T_a}{R_{ao}}+\frac{T_o-T_e}{R_{eo}}+Q_c+Q_r.
\]
The exact topology and three-way routing are \ScaleBridge{} ratified choices.  Detailed multi-zone RC literature demonstrates the broader use of room/wall thermal nodes and flexible RC structures \citep{BelicSliskovicHocenski2021MultiZoneRC,VallianosAthienitisDelcroix2022AutomaticMultiZoneRC}.

\section{Default simulation initialization}
For the default Phase-E simulation policy, observed air temperature initializes the observed state and every latent thermal state in the same zone:
\[
T_\ell^{(z)}(t_0)=T_a^{(z)}(t_0)\quad\text{for every latent state }\ell.
\]
Thus 2R2C/3R2C use $T_m(t_0)=T_a(t_0)$ and 4R3C uses $T_e(t_0)=T_m(t_0)=T_a(t_0)$. This is an initialization assumption, not a claim of equality after integration begins.

\bibliographystyle{plainnat}
\bibliography{../../references/ScaleBridge_Thermal_Modeling_References}
\end{document}
